\section{Introduction}
\label{sec:intro}

A wireless sensor node spends most of its energy on its radio, and in nearly all
deployments its battery cannot be replaced~\cite{akyildiz2002survey}. Lifetime
is therefore decided by how the network organizes its transmissions rather than
by any property of the sensing hardware. Clustering is the standard response:
nodes elect a subset of themselves as cluster heads, members send short-range
packets to their head, and the head fuses them into one packet and pays the
single expensive transmission to the base station. Because that role is
expensive it has to rotate, and the question every clustering protocol answers
is which nodes should carry it this round.

Twenty-five years of answers fall into three broad generations. The first is
heuristic. LEACH rotates the role by an independent probabilistic threshold with
an epoch constraint~\cite{heinzelman2000leach}; PEGASIS abandons clusters
entirely for a greedy chain in which each node talks only to its nearest
neighbor~\cite{lindsey2002pegasis}; TEEN and its extension APTEEN leave the
topology alone and instead suppress transmissions that fall below a sensed-value
threshold~\cite{manjeshwar2001teen,manjeshwar2002apteen}. Later heuristics
such as HEED refine the election rule without changing its
character~\cite{younis2004heed,abbasi2007clustering}. The second generation
replaces the heuristic with an explicit objective. Head selection is posed as a
multi-objective optimization problem and solved with
NSGA-II~\cite{deb2002nsga2}, or encoded in a type-2 fuzzy rule base that reasons
over residual energy, distance and local density under linguistic
uncertainty~\cite{mendel2002type2,karnik2001centroid}. The third generation
learns the selection function instead of specifying
it~\cite{alsheikh2014ml}: a self-organizing map places heads by
unsupervised vector quantization~\cite{kohonen1990som}, a deep Q-network treats
rotation as a sequential decision problem~\cite{mnih2015dqn,sutton2018rl}, and
a graph convolutional network scores candidates from the topology's adjacency
structure~\cite{kipf2017gcn,wu2021gnnsurvey}.

The generations are usually compared by reading their published tables side by
side, and that comparison does not hold. Each protocol was proposed against a
different simulator, field size, node count, packet length and energy
budget~\cite{kurkowski2005incredibles,pawlikowski2002credibility,andel2006credibility},
almost always
against LEACH alone, and almost always on a lossless channel. More
consequentially, the assumptions that decide the outcome are rarely stated.
Whether a centralized protocol is charged for the uplink that tells the base
station each node's residual energy can be worth a quarter of the entire energy
budget over a long run. Whether fusion is lossless in bits determines how much a
chain protocol gains over a clustered one. Whether lifetime is reported at first
death, half death or last death can reverse a ranking outright, and the author
chooses which to report. A margin obtained under one set of these choices is not
evidence about the algorithm; it is partly evidence about the harness.

This paper measures the three generations in one harness. Nine protocols
spanning all of them, plus a direct-transmission baseline, run on a single
engine over 30 paired trials in each of two channel configurations, for 600 runs
in the headline study, and over a further 1{,}350 runs in a node-count by
field-area sweep. Run index $i$ seeds the topology, the per-link shadowing and
the entire sensed-value stream identically for every protocol, so a comparison
at a given index is genuinely paired rather than an average over unrelated
worlds. Fairness is enforced structurally instead of by convention: the engine
owns energy accounting through a single charging function, packet sizing, the
aggregation charge, the channel and its ARQ retries, and the counting of
delivered readings. A protocol returns only a cluster structure and a route
shape, and a static audit verifies that no protocol implementation writes to the
energy or liveness arrays. Every protocol runs on the same numerical substrate,
hand-written NumPy including the backpropagation for both learned models,
because per-round setup time is a reported metric and mixing frameworks would
corrupt it. No protocol is tuned to reproduce any published number;
disagreement with the source tables is expected and is not treated as error.

We also state up front the one assumption that most favors part of the field.
Centralized protocols are not charged for the uplink that carries node state to
the sink, because that state is treated as piggybacked on data traffic the
network already sends. This is the standard convention in the centralized
clustering literature, but it is still a choice, and Section~\ref{sec:validity}
measures what it is worth: LEACH spends 12.73\% of its consumed energy on
control traffic against 2.33\% to 3.95\% for the centralized five, so every
comparison that crosses those classes carries a subsidy of roughly nine to ten
percentage points of the budget.

Five results are worth stating before the details.

Transmit power is not a background parameter but a determinant of the ranking.
Repeating the entire study on a lossless channel removes the first-death
advantage of the type-2 fuzzy protocol and of the DQN over LEACH altogether:
$+417$ and $+426$ rounds at $p_{\mathrm{Holm}} = 0.00045$ under packet loss,
falling to $+17$ and $+22$ rounds and no significance once loss is removed.

Direct measurement of where the energy goes explains why, and the explanation is
not the obvious one. Packet error is effectively zero below 120~m, and every
protocol's \emph{mean} head-to-sink distance sits below that, so the mean cannot
be the cause. The tail can. LEACH puts 29.9\% of its head-rounds beyond 120~m
and spends 6.39\% of its budget on retransmissions; the DQN puts 5.1\% of its
head-rounds out there and spends 2.31\%. Across the six single-hop protocols the
share of head-rounds beyond 120~m predicts retry energy share with $r = 0.95$.
What the newer protocols demonstrably do well is keep the far tail of the
head-to-sink hop out of the error waterfall, and that skill pays only in
proportion to how lossy the hop is.

The effect is specific to first death. Differences in area under the alive-node
curve and in readings delivered hold at $p_{\mathrm{Holm}} = 0.00045$ in both
channels.

No single lifetime figure summarizes clustering, because it buys a later first
death by bringing the last death forward. LEACH reaches first death at 1{,}046
rounds against the baseline's 114, a factor of 9.2, while dying out completely
at 2{,}191 rounds against the baseline's 3{,}202. The GCN is the extreme case,
with a first death 3.1 times earlier than LEACH's and a last death 45\% later.

Finally, the advantage of clustering is a property of an operating regime rather
than of clustering as such. Across a $3\times3$ grid of node counts and field
areas, tripling the node count moves first death by a factor of 0.88 to 1.26
while tripling the field side moves it by a factor of 5 to 7. In the
$50\times50$~m field, where no link reaches the error waterfall, LEACH is
outlived by the no-clustering baseline.

The reactive protocols sit outside all of this because they survive by not
reporting: TEEN and APTEEN fail to reach last death within the 7{,}000-round cap
in all 30 trials, at a data yield of 0.11 and 0.12 against roughly 0.9 for the
proactive protocols.

The contributions are as follows. First, a single simulation engine in which
nine clustering protocols from three methodological generations execute under
provably identical physical conditions, with fairness enforced by construction
rather than by author discipline. Second, a paired experimental design, common
topology, shadowing and sensed-data stream per trial index, analyzed with paired
sign-flip permutation tests, paired bootstrap intervals and Holm correction
across the protocol family, in place of the unpaired mean-and-deviation
comparison the literature normally reports. Third, a channel sensitivity pair
run as a rule rather than as an appendix: a conclusion that survives both the
lossy and the ideal configuration is a statement about the protocol, and one
that flips is a statement about transmit power. Applying the rule retracts a
comparison that the lossy channel alone would have supported. Fourth, a
node-count by field-area sweep that identifies the regime boundary from an
independent direction and tests whether the frozen learned policies transfer off
their training distribution. Fifth, mechanistic explanations obtained by
measurement rather than argument, including the energy split by category, the
distribution of cluster-head distance relative to the radio crossover, and
cluster-head rotation concentration. Sixth, a full disclosure of the modeling
choices that bound these results, ordered by how far each could move a
conclusion, together with the reproduction artifacts.

Section~\ref{sec:model} describes the system and radio model,
Section~\ref{sec:protocols} the nine protocols and their fairness-relevant
implementation details, Section~\ref{sec:design} the experimental design and
statistical procedure, Section~\ref{sec:results} the results,
Section~\ref{sec:validity} the threats to validity, and
Section~\ref{sec:conclusion} concludes.
